
\begin{comment}
\section{Some examples of convex biproduct categories}\label{app:examples}

In this appendix we illustrare several examples of convex biproduct categories:  the Kleisli category for the subdistribution monad,  its continuous version and, finally, a result about finitely partially additive categories.



\subsection{The Kleisli category of the subdistribution monad}

\begin{proposition}\label{prop:Klconvexbiproduct} %\marginpar{Aggiungi dimostrazione}
$\KlD$ is a convex biproduct category
\end{proposition}
\begin{proof}%[Proof of Proposition \ref{prop:Klconvexbiproduct}]
 Coproducts are induced by the coproducts in $\Sets$, and the empty set $\zero$ is both initial and final. Indeed, for every object $X$, the unique arrow $X\to \zero$ sends an element $x\in X$ into the null subdistribution (which is the unique element of $\Dis(\zero)$). 
    
$\KlD$ is $\Cat{PCA}$-enriched since, for every $X,Y\in \KlD$, the hom-set $\KlD[X,Y]$ is a pointed convex algebra with the operations defined as follows: $+_p(f,g)(y|x)\defeq p\cdot f(y|x) + (1-p)\cdot g(y|x)$ for all $f,g\colon X \to Y$ and $p\in (0,1)$, and $\star_{X,Y}(y|x) \defeq 0$ (i.e., it is the arrow which sends every $x$ into the null subdistribution over $Y$). An easy calculation shows that arrow composition is a morphism of pointed convex algebras, i.e., it preserves the operations $+_p$ and $\star_{X,Y}$. For instance $(f+_pg) ;h= (f;h +_p g;h)$ for all $f,g\colon X \to Y$ and $h\colon Y \to Z$ is obtained as follows:
\begin{align}
    ((f+_pg) ;h)(z|x) &= \sum_{y\in Y} (f+_pg)(y|x) \cdot h(z|y) \tag*{} \\
    &= \sum_{y\in Y} (p\cdot f(y|x) + (1-p)\cdot g(y|x)) \cdot h(z|y) \tag*{}\\
    &= p\cdot \sum_{y\in Y} f(y|x) \cdot h(z|y) + (1-p)\cdot \sum_{y\in Y} g(y|x) \cdot h(z|y) \tag*{}\\
    &= p\cdot (f;h)(z|x) + (1-p)\cdot (g;h)(z|x) \tag*{}\\
    &= (f;h +_p g;h)(z|x)\text{.}\tag*{}
    \end{align}
Now observe that $X_1\oplus X_2$ with the projections $\pi_i\colon X_1 \oplus X_2 \to X_i$ for $i=1,2$ given by 
\[ \pi_1(x_1|z)\defeq \begin{cases}
    \delta_{x_1} & z=\iota_1(x_1)\\
    0 & otherwise
\end{cases} \qquad \pi_2(x_2|z)\defeq \begin{cases}
     \delta_{x_2} & z=\iota_2(x_2)\\
        0 & otherwise
\end{cases}\]
 is a convex product of $X_1$ and $X_2$. Indeed, given arrows $f_1\colon A\to X_1$ and $f_2\colon A\to X_2$ and $p_1,p_2 \in (0,1)$ such that $p_1+p_2\le 1$, there exists an arrow $\langle f_1,f_2\rangle_{p_1,p_2}\colon A \to X_1\oplus X_2$ defined as follows:
\[ \langle f_1,f_2\rangle_{p_1,p_2}(z|a) \defeq \begin{cases}
    p_1\cdot f_1(x_1|a) & z=\iota_1(x_1)\\
    p_2\cdot f_2(x_2|a) & z=\iota_2(x_2) \text{.}
\end{cases}\]
Observe that, for $i\in \{1,2\}$,
\begin{align*}
\langle f_1,f_2\rangle_{p_1,p_2} ; \pi_i (x_i | a) = & \sum_{z\in X_1+X_2} \langle f_1,f_2\rangle_{p_1,p_2}(z|a) \cdot \pi_i (x_i | z) \\
= & \langle f_1,f_2\rangle_{p_1,p_2} (\iota_i(x_i)|a)\\
= & p_i\cdot f_i (x_i|a)
\end{align*}
namely, it holds that $(\langle f_1,f_2\rangle_{p_1,p_2} );\pi_i = p_i\cdot f_i$. This is the unique arrow with such property since, for any other arrow $h'\colon A \to X_1\oplus X_2$ such that $h';\pi_i = p_i\cdot f_i$ for all $i=1,2$ it follows that 
\begin{align*}
  h'(\iota_1(x_1)|a) &=    h'(\iota_1(x_1)|a) \cdot 1\\
    & =\sum_{z\in X_1\oplus X_2} h'(z|a) \cdot \pi_1(x_1|z) \tag*{}\\
   &= h' ; \pi_1(x_1|a) \\
    &= p_1\cdot f_1(\iota_1(x_1)|a) \tag*{}
\end{align*}
and similarly for $\pi_2$ one obtains that $h'(\iota_2(x_2)|a) = p_2\cdot f_2(\iota_2(x_2)|a)$, i.e., $h' = \langle f_1,f_2\rangle_{p_1,p_2}$. 
Simple computations confirms that \eqref{eq:delta} holds.
\end{proof}



 %   \marginpar{Nuova prova per il Thm 23.}
 

%\marginpar{mi sono accorto che manca la proposizione che dice che i prodotti n-ary sono preservati. Potremmo aggiungere un corollario dopo la Proposizione 30, mi sembra si dimostri allo stesso identico modo. Ho aggiunto il riferimento a questa prop anche nella prova del teorema 12.}



\subsection{Substochastic Markov kernels on standard Borel spaces}
\end{comment}


\section{Substochastic Markov kernels on standard Borel spaces}\label{app:examples}

Consider the category of \textit{standard Borel spaces} and measurable functions denoted with $\mathbf{BorelMeas}$.
Given a {standard Borel space} $(X,\Sigma_X)$, denote with $\mathcal{G}_{\le}(X)$ the set of subprobability measures on $X$, i.e. measurable functions $\mu:(X,\Sigma_X)\to ([0,1],\mathcal{B}([0,1]))$ such that $\mu(X)\le 1$, where $\mathcal{B}([0,1])$ is the Borel $\sigma$-algebra on $[0,1]$.  This assignment extends to a functor $\mathcal{G}_{\le}:\mathbf{BorelMeas}\to \mathbf{BorelMeas}$ which corresponds to the subdistribution version of the \textit{Giry monad}~\cite{giry2006categorical}. The functor $\mathcal{G}_{\le}$ is a symmetric monoidal monad and its Kleisli category, $\Kl({\mathcal{G}_{\le}})$, is the category of standard Borel spaces and substochastic Markov kernels. A substochastic Markov kernel from $(X,\Sigma_X)$ to $(Y,\Sigma_Y)$ is a function $f\colon \Sigma_Y\times X\to [0,1]$ such that for all $x\in X$, $f(-|x)\colon \Sigma_Y\to [0,1]$ is a subprobability measure on $Y$ and for all $U\in \Sigma_Y$, $f(U|-)\colon X\to [0,1]$ is a measurable function.
The identity arrow on $(X,\Sigma_X)$ is the Markov kernel $\delta_X$ such that for all $x\in X$ and $U\in \Sigma_X$, $\delta_X(U|x)=1$ if $x\in U$ and $0$ otherwise. The composition of arrows in $\Kl({\mathcal{G}_{\le}})$ is given by the Chapman-Kolmogorov equation: for all  substochastic Markov kernels $f\colon (X,\Sigma_X)\to (Y,\Sigma_Y)$ and $g\colon (Y,\Sigma_Y)\to (Z,\Sigma_Z)$, $f;g\colon (X,\Sigma_X)\to (Z,\Sigma_Z)$ is defined as the Markov kernel such that for all $x\in X$ and $U\in \Sigma_Z$,
\begin{equation*}
(f;g)(U|x)\coloneqq \int_Y g(U|y)\, f(dy|x)
\end{equation*}
   
\begin{proposition}
    \label{prop: SubStoch is convex biproduct}
    $\Kl({\mathcal{G}_{\le}})$ is a convex biproduct category.
\end{proposition}
\begin{proof}
Coproducts in $\Kl({\mathcal{G}_{\le}})$ are given by the coproducts in $\mathbf{BorelMeas}$, i.e. the disjoint union of standard Borel spaces. The initial and terminal object is the empty standard Borel space. The $\mathbf{PCA}$-enrichment is given by the pointwise convex combination of Markov kernels: for all Markov kernels $f,g\colon (X,\Sigma_X)\to (Y,\Sigma_Y)$ and $p\in [0,1]$, $f +_p g\colon (X,\Sigma_X)\to (Y,\Sigma_Y)$ is defined  for all $x\in X$ and $U\in \Sigma_Y$ as  $f+_p g(U|x)\coloneqq p\cdot f(U|x)+(1-p)\cdot g(U|x)$. The zero arrow $\star_{(X,\Sigma_X),(Y,\Sigma_Y)}\colon (X,\Sigma_X)\to (Y,\Sigma_Y)$ is the Markov kernel such that for all $x\in X$ and $U\in \Sigma_Y$, $\star_{(X,\Sigma_X),(Y,\Sigma_Y)}(U|x)=0$. The axioms of $\mathbf{PCA}$ are satisfied since they are so in $[0,1]$. A simple computation shows that $+_p$ and $\star$ are compatible with composition. For instance, %as done in Proposition \ref{prop:Klconvexbiproduct},
 we obtain that $(f+_p g);h=f;h +_p g;h$ for all Markov kernels $f,g\colon (X,\Sigma_X)\to (Y,\Sigma_Y)$, $h\colon (Y,\Sigma_Y)\to (Z,\Sigma_Z)$ and $p\in [0,1]$ in the following way
\begin{align}
\big((f+_{p}g);h\big)(C\mid x)
&= \int_{Y} h(C\mid y)\,(f+_{p}g)(dy\mid x)
\tag*{} \\
&= \int_{Y} h(C\mid y)\,\big(p\,f(dy\mid x)+(1-p)\,g(dy\mid x)\big)
\tag*{} \\
&= p\int_{Y} h(C\mid y)\,f(dy\mid x) \;+\; (1-p)\int_{Y} h(C\mid y)\,g(dy\mid x)
\tag{linearity of the integral}\\
&= p\,(f;h)(C\mid x) \;+\; (1-p)\,(g;h)(C\mid x)
\tag*{} \\
&= \big((f;h)+_{p}(g;h)\big)(C\mid x).
\tag*{}
\end{align}






Now we prove that $(X_1,\Sigma_{X_1})\overset{\pi_1}{\leftarrow}(X_1,\Sigma_{X_1})\oplus (X_2,\Sigma_{X_2})\overset{\pi_2}{\rightarrow} (X_2,\Sigma_{X_2})$ is a binary convex product in $\Kl({\mathcal{G}_{\le}})$. The Markov kernels $\pi_1$ and $\pi_2$ are defined as follows: for all $z\in X_1\oplus X_2$ and $U\in \Sigma_{X_1}$, $\pi_1(U|z)=1$ if $z\in \iota_1(U)$ and $0$ otherwise; for all $z\in X\oplus Y$ and $V\in \Sigma_{X_2}$, $\pi_2(V|z)=1$ if $z\in \iota_2(V)$ and $0$ otherwise. Given $f\colon (A,\Sigma_A)\to (X_1,\Sigma_{X_1})$, $g\colon (A,\Sigma_A)\to (X_2,\Sigma_{X_2})$ and $p_1,p_2\in [0,1]$ such that $p_1+p_2\leq 1$, we prove that there is a unique Markov kernel $h\colon (A,\Sigma_A)\to (X_1,\Sigma_{X_1})\oplus (X_2,\Sigma_{X_2})$ such that $h;\pi_1 = p_1\cdot f$ and $h;\pi_2 = p_2\cdot g$. Existence is provided by the Markov kernel $h$ defined as follows: for all $a\in A$, $U\in \Sigma_{X_1}$ and $V\in \Sigma_{X_2}$,
\[ h(\iota_1(U)|a)\coloneqq p_1\cdot f(U|a) \quad h(\iota_2(V)|a)\coloneqq p_2\cdot g(V|a) \]
where for a generic element of $\Sigma_{X_1\oplus X_2}$ $h$ is extended by countable additivity. Uniqueness is proved as follows: suppose that there is also another $h'$ with the same property. Then, for all $a\in A$, $U\in \Sigma_{X_1}$ and $V\in \Sigma_{X_2}$, 
\begin{align}
    h'(\iota_1(U)\mid a)
    &= \int_{X_1 \oplus X_2} \mathbf{1}_{\iota_1(U)}(z)\, h'(dz\mid a)
    \tag{def. of integration w.r.t.\ $h'(\cdot\mid a)$}\\
    &= \int_{X_1 \oplus X_2} \pi_1(U\mid z)\, h'(dz\mid a)
    \tag{def. of indicator function}\\
    &= (h';\pi_1)(U\mid a)
    \tag{kernel composition}\\
    &= (p_1\cdot f)(U\mid a)
    \tag{hp. $h';\pi_1 = p_1\cdot f$}.
\end{align}


Similarly, $h'(\iota_2(V)|a)=(p_2\cdot g)(V|a)$. Since every element of $\Sigma_{X_1\oplus X_2}$ is obtained by countable unions of elements of the form $\iota_1(U)$ and $\iota_2(V)$, $h'$ is completely determined by the above equalities, hence $h'=h$.
\end{proof}

%%%%%%%COMMENTO TUTTA A ROBA SU FIN PAC; ORA NEL TESTO PRINCIPALE%%%%%%%%%%%%%%
\begin{comment}

\subsection{On finitely partially additive categories and convex biproduct categories}\label{app:eff}
Our probabilistic tape diagrams seem to be closely related to effectuses \cite{introductioneffectus}. These are founded on finitely partially additive categories \cite{chophd} that, as we illustrate in Proposition \ref{prop:fpac}, are  convex biproduct categories under suitable conditions.

We refer the reader to \cite[Definition 3.1.2]{chophd} for the definition of finitely partially additive categories. For the purpose of this appendix, it is enough to recall that 
while convex bicategories are enriched over pcas, finitely partially additive categories are enriched over partial commutative monoids.
A \emph{partial commutative monoid}, pcm for short, consists of a set $X$, together with a partial function $\circledvee\colon X \times X \to X$ and an element $0\in X$ satisfying the following laws, where we write $x\bot y$ to means that $x\circledvee y$ is defined:
\begin{itemize}
\item if $x\bot y$ and $x\circledvee y \bot z$, then $y \bot z$, $x \bot y\circledvee z$ and $(x\circledvee y) \circledvee z = x\circledvee (y \circledvee z)$;
\item if $x\bot y$, then $y\bot x$ and $x \circledvee y = y \circledvee x$;
\item $0\bot x$ and $0\circledvee x =x$. 
\end{itemize}
%An homomorphism of pcms is a function $f\colon X \to Y$ such that
%\begin{itemize}
%\item if $x \bot y$, then $f(x)\bot f(y)$ and $f(x\circledvee y) = f(x) \circledvee f(y)$;
%\item $f(0)=0$.
%\end{itemize}
%We write $\Cat{PCM}$ for the category of pcms and their morphisms and $\Cat{PCMCat}$ for the category of categories enriched over $\Cat{PCM}$.


Our starting observation is that any pca $(X,+_p, \star)$ carries the structure of a pcm $(X,\circledvee, 0)$ where $0\defeq \star$ and for all $x,y\in X$,
\begin{equation}\label{eq:defbot}
x\bot y \text{ iff } \exists x',y'\in X, \, p,q\in[0,1] ,\; x=p\cdot x', \; y=q\cdot y', \; p+q\leq 1
\end{equation}
and in this case
\begin{equation}\label{eq:defcircledvee}
x \circledvee y \defeq x'+_p(y'+_{\frac{q}{1-p}}\star)\text{.}
\end{equation}
%The operation $\circledvee$ is well defined and indeed $(X,\circledvee, 0)$ is a pcm as shown in the following result.
\begin{lemma}\label{lemma:pca-pcm}
Any pca gives rise to a pcm with the structure described above.
\end{lemma}
\begin{proof}
Given any pca $(X,+_p,\star)$, consider the structure $(X,\circledvee,0)$ defined as above. We first prove that $\circledvee$ is well defined, then we prove that the axioms of pcm are satisfied. Assume that $x\bot y$, which means that there exist $x',y'\in X$ and $p,q\in [0,1]$ such that $x=p\cdot x'$, $y=q\cdot y'$ and $p+q\leq 1$. Assume also that $x\bot y$ via $x'',y''\in X$ and $r,s\in [0,1]$ such that $x=r\cdot x''$, $y=s\cdot y''$ and $r+s\leq 1$. We have to prove the equality
\begin{equation}\label{eq:pcmbendefinita}
    x'+_p(y'+_{\frac{q}{1-p}}\star) = x''+_r(y''+_{\frac{s}{1-r}}\star)\text{.}\end{equation}
In order to do that, we assume without loss of generality that $p\leq r$ and that $s\leq q$ (the other cases can be treated similarly). Now we show that both sides of \eqref{eq:pcmbendefinita} are equal to $x'+_p(y''+_{\frac{s}{1-p}}\star)$. We have
\begin{align*}
    x'+_p(y'+_{\frac{q}{1-p}}\star) & = x'+_p((\frac{q}{1-p}\cdot y')) \tag{\ref{eq:def somma n pca}}\\
    & = x'+_p((\frac{s}{1-p}\cdot y'')) \tag{$q\cdot y'=s\cdot y''$}\\
    & = x'+_p(y''+_{\frac{s}{1-p}}\star) \tag{\ref{eq:def somma n pca}}
    \end{align*}
    Similarly,
    \begin{align*}
 x''+_r(y''+_{\frac{s}{1-r}}\star) & = y''+_s (x''+_{\frac{r}{1-s}}\star) \tag{using axioms in \ref{eq:pca}}\\
    & = y''+_s(\frac{r}{1-s}\cdot x'') \tag{\ref{eq:def somma n pca}}\\
    & = y''+_s(\frac{p}{1-s}\cdot x') \tag{$p\cdot x'=r\cdot x''$}\\
    & = y''+_s(x'+_{\frac{p}{1-s}}\star) \tag{\ref{eq:def somma n pca}}\\
    & = x'+_p(y''+_{\frac{s}{1-p}}\star ) \tag{using axioms in \ref{eq:pca}}
\end{align*}
Hence the equality in \eqref{eq:pcmbendefinita} holds and $\circledvee$ is well defined. Now we prove the axioms of pcm. We illustrate the proof of symmetry, the other axioms can be proved similarly. Assume that $x\bot y$, and hence there exist $x',y'\in X$ and $p,q\in [0,1]$ such that $x=p\cdot x'$, $y=q\cdot y'$, $p+q\leq 1$. Then $y\bot x$ via $y',x'$ and $q,p$. 
Moreover, 
\begin{align*}
    x\circledvee y & = x'+_p(y'+_{\frac{q}{1-p}}\star) \tag{\eqref{eq:defcircledvee}}\\
    & = (y'+_{\frac{q}{1-p}}\star) +_{1-p} x' \tag{symmentry in \eqref{eq:pca}}\\
    & = y'+_q(\star +_{\frac{1-p-q}{1-q}} x') \tag{associativity in \eqref{eq:pca}}\\
    & = y'+_q(x'+_{\frac{p}{1-q}}\star) \tag{symmetry in \eqref{eq:pca}}\\
& = y\circledvee x \tag{\eqref{eq:defcircledvee}}
\end{align*}
\end{proof}


\begin{proposition}
Any pca-enriched category is a pcm-enriched category with the enrichment given by the pcm described above.
\end{proposition}
\begin{proof}
Let $\Cat{C}$ be a pca-enriched category. By Lemma \ref{lemma:pca-pcm}, any homset $\Cat{C}[X,Y]$ is a pcm. To conclude that it is $\Cat{C}$ is pcm-enriched is enough to prove  for all properly typed arrows
\begin{itemize}
\item if $g\bot h$, then $f;(g \circledvee h) = (f;g) \circledvee (h;h) $;
\item if $f\bot g$, then  $(f \circledvee g); h = (f;h) \circledvee (g;h)$;
\item $f;0=0=0;f$.
\end{itemize}
We illustrate below the proof of $(f \circledvee g); h = (f;h) \circledvee (g;h)$.

Assume that $f\bot g$, then by \eqref{eq:defbot}, there exist arrows $f',g'$ and  $p,q\in [0,1]$ such that
\[f=p\cdot f' \quad g=q\cdot g' \qquad p+q\leq 1\text{.}\]
%Thus, by \eqref{}, \marginpar{INSERT DEFINITION ABOVE} $f\circledvee g = f' +_{p} (g'+_{\frac{q}{1-p}} \star)$. 
Thus 
\begin{align*}
(f\circledvee g); h & = ((f' +_{p} (g'+_{\frac{q}{1-p}} \star));h \tag{\eqref{eq:defcircledvee}} \\
&  = ((f';h +_{p} (g';h +_{\frac{q}{1-p}} \star)) \tag{$\Cat{C}$ is pca-enriched} \\
&= f;h \circledvee g;h
\end{align*}
where the last step holds by \eqref{eq:defcircledvee} and the fact that,  by Lemma \ref{lemma:initialproperties2}.\ref{lemma:p;}, $p\cdot (f';h) = (p\cdot f'); h = f;h$ and $q\cdot (g';h) = (q\cdot g'); h = g;h$.
\end{proof}



\begin{proposition}\label{prop:fpac}
Let $\Cat{C}$ be a pca-enriched category. If $\Cat{C}$ is a finitely partially additive category, then it is a convex biproduct category.
\end{proposition}
\begin{proof}
Since $\Cat{C}$ is a finitely partially additive category, then by definition (\cite[Definition 3.1.2]{chophd} ) it has finite coproducts and zero morphisms. Thus, it has zero object, i.e., initial and final object coincide, and the equality in \eqref{eq:delta} holds.

In order to show that $\oplus$ is a convex product in $\Cat{C}$, take two arrows $f\colon A\to X_1$ and $g\colon A \to X_2$ and $p_1,p_2\in (0,1)$ such that $p_1+p_2\leq 1$.
Consider now the arrows $(p_1\cdot f) ; \iota_1$ and $(p_2 \cdot g) ;\iota_2$ in the pcm  $\Cat{C}[A , X_1\oplus X_2]$. Since by Lemma \ref{lemma:initialproperties2}.\ref{lemma:p;}  $(p_1\cdot f) ; \iota_1 = p_1 \cdot(f;\iota_1)$ and 
$(p_2\cdot g) ; \iota_2 = p_2 \cdot(g;\iota_2)$ and since $p_1+p_2\leq 1$, then by definition of $\bot$ in \eqref{eq:defbot}, it holds that 
\[(p_1\cdot f) ; \iota_1 \; \bot \; (p_2 \cdot g ) ;\iota_2 \text{.}\]
We thus can take $\langle f,g\rangle_{p_1,p_2}$ as $((p_1\cdot f) ; \iota_1)  \circledvee ((p_2 \cdot g ) ;\iota_2)\colon A \to X_1\oplus X_2$ which, by  \cite[Proposition 3.1.9]{chophd}, is the unique arrow such that \[\langle f,g\rangle_{p_1,p_2} ; \pi_1 = p_1\cdot f \quad \text{ and } \quad \langle f,g\rangle_{p_1,p_2} ; \pi_2 = p_2\cdot g\text{.}\]
%The compatibility condition in \eqref{eq: assioma aggiuntivo} follows immediately from \cite[Proposition 3.1.8]{chophd}.
\end{proof}


    

    

\begin{remark}\label{rmk:eff}
The reader may now wonder whether any convex biproduct category is also a finitely partially additive category. This is not the case in general. 

The definition of finitely partially additive category requires that \emph{compatible maps}, i.e., those $f_1\colon A \to X$, $f_2\colon A \to X$ for which there exist some $h\colon A\to X \oplus X$ such that $h;\pi_i=f_i$ (see \cite[Definition 3.1.1]{chophd} for more details), are summable. In a convex biproduct category there might exist compatible $f_1$ and $f_2$ that are not summable,  namely there are no $f_1',f_2'$ such that $f_1=p_1\cdot f_1'$ and $f_2=p_2\cdot f_2'$ for some scalar $p_1+p_2 \leq 1$. 

The essential difference can be understood as follow: summability in convex biproduct categories is fixed a-priori by the structure of pca, while in partially additive category is left to the structural properties of morphisms. For instance, while in convex biproduct categories $\oplus$ behaves as  convex product, in partially additive category $\oplus$ acts as a \emph{partial} biproduct: the product property only holds for compatible maps.   

This difference is crucial for our approach which is grounded on monoidal algebra. Indeed, in an arbitrary finitely partially additive category, one cannot define $\diagp{X}\colon X \to X\oplus X$ unless some notion of scalar is explicitly required. 

Scalars appears in \emph{effectuses} \cite{introductioneffectus} as arrows of type $1 \to 1$ where $1$ is a distinguished object. Interestingly in $\CatT{\Cat{C}}$, for some monoidal category $(\Cat{C}, \otimes, 1)$,
every  $p\in [0,1]$ defines the arrow $\diagp{1}; (\id{1} \oplus \, \bang{1}) \colon 1 \to 1$. Nevertheless, not all the arrows in $\CatT{\Cat{C}}[1,1]$ are $p\in[0,1]$ (unless $\Cat{C}[1,1]$ is the singleton set): see for instance the syntactic categorty $\Cat{T}(\Diag{B})$ in Section \ref{sec:probboolcircuits}. To make the correspondence with effectuses more precise, one needs to consider more structure for tapes, such as their \emph{copy-discard} variants introduced in \cite{bonchi2025tapediagramsmonoidalmonads}. 
\end{remark}



%%%%%%%%%%SOTTO LA VERSIONE NUOVA DELLA SEZIONE FINPAC CHE HO RISCRITTO%%%%%%%% 

%\begin{comment}

\subsection{Nuova corrispondenza FinPAc e CBC}
\marginpar{Qui metto una riformulazione della sezione precedente, alla luce del fatto che la transitività non mi torna nel passaggio PCA--->PCM}

Our probabilistic tape diagrams seem to be closely related to effectuses \cite{introductioneffectus}. These are founded on finitely partially additive categories \cite{chophd} that, as we illustrate in Proposition \ref{prop:fpac}, are  convex biproduct categories under suitable conditions.

We refer the reader to \cite[Definition 3.1.2]{chophd} for the definition of finitely partially additive categories. For the purpose of this appendix, it is enough to recall that 
while convex bicategories are enriched over pcas, finitely partially additive categories are enriched over partial commutative monoids.
A \emph{partial commutative monoid}, pcm for short, consists of a set $X$, together with a partial function $\circledvee\colon X \times X \to X$ and an element $0\in X$ satisfying the following laws, where we write $x\bot y$ to means that $x\circledvee y$ is defined:
\begin{itemize}
\item if $x\bot y$ and $x\circledvee y \bot z$, then $y \bot z$, $x \bot y\circledvee z$ and $(x\circledvee y) \circledvee z = x\circledvee (y \circledvee z)$;
\item if $x\bot y$, then $y\bot x$ and $x \circledvee y = y \circledvee x$;
\item $0\bot x$ and $0\circledvee x =x$. 
\end{itemize}
%An homomorphism of pcms is a function $f\colon X \to Y$ such that
%\begin{itemize}
%\item if $x \bot y$, then $f(x)\bot f(y)$ and $f(x\circledvee y) = f(x) \circledvee f(y)$;
%\item $f(0)=0$.
%\end{itemize}
%We write $\Cat{PCM}$ for the category of pcms and their morphisms and $\Cat{PCMCat}$ for the category of categories enriched over $\Cat{PCM}$.


Our starting observation is that any pca $(X,+_p, \star)$ carries the structure of a pcm $(X,\circledvee, 0)$ where $0\defeq \star$ and for all $x,y\in X$,
\begin{equation}\label{eq:defbot}
x\bot y \text{ iff } \exists x',y'\in X, \, p,q\in[0,1] ,\; x=p\cdot x', \; y=q\cdot y', \; p+q\leq 1
\end{equation}
and in this case
\begin{equation}\label{eq:defcircledvee}
x \circledvee y \defeq x'+_p(y'+_{\frac{q}{1-p}}\star)\text{.}
\end{equation}
%The operation $\circledvee$ is well defined and indeed $(X,\circledvee, 0)$ is a pcm as shown in the following result.
\begin{lemma}%\label{lemma:pca-pcm}
Any free pca gives rise to a pcm with the structure described above.
\end{lemma}
\begin{proof}
Given any pca $(X,+_p,\star)$, consider the structure $(X,\circledvee,0)$ defined as above. We first prove that $\circledvee$ is well defined, then we prove that the axioms of pcm are satisfied. Assume that $x\bot y$, which means that there exist $x',y'\in X$ and $p,q\in [0,1]$ such that $x=p\cdot x'$, $y=q\cdot y'$ and $p+q\leq 1$. Assume also that $x\bot y$ via $x'',y''\in X$ and $r,s\in [0,1]$ such that $x=r\cdot x''$, $y=s\cdot y''$ and $r+s\leq 1$. We have to prove the equality
\begin{equation}\label{eq:pcmbendefinita}
    x'+_p(y'+_{\frac{q}{1-p}}\star) = x''+_r(y''+_{\frac{s}{1-r}}\star)\text{.}\end{equation}
In order to do that, we assume without loss of generality that $p\leq r$ and that $s\leq q$ (the other cases can be treated similarly). Now we show that both sides of \eqref{eq:pcmbendefinita} are equal to $x'+_p(y''+_{\frac{s}{1-p}}\star)$. We have
\begin{align*}
    x'+_p(y'+_{\frac{q}{1-p}}\star) & = x'+_p((\frac{q}{1-p}\cdot y')) \tag{\ref{eq:def somma n pca}}\\
    & = x'+_p((\frac{s}{1-p}\cdot y'')) \tag{$q\cdot y'=s\cdot y''$}\\
    & = x'+_p(y''+_{\frac{s}{1-p}}\star) \tag{\ref{eq:def somma n pca}}
    \end{align*}
    Similarly,
    \begin{align*}
 x''+_r(y''+_{\frac{s}{1-r}}\star) & = y''+_s (x''+_{\frac{r}{1-s}}\star) \tag{using axioms in \ref{eq:pca}}\\
    & = y''+_s(\frac{r}{1-s}\cdot x'') \tag{\ref{eq:def somma n pca}}\\
    & = y''+_s(\frac{p}{1-s}\cdot x') \tag{$p\cdot x'=r\cdot x''$}\\
    & = y''+_s(x'+_{\frac{p}{1-s}}\star) \tag{\ref{eq:def somma n pca}}\\
    & = x'+_p(y''+_{\frac{s}{1-p}}\star ) \tag{using axioms in \ref{eq:pca}}
\end{align*}
Hence the equality in \eqref{eq:pcmbendefinita} holds and $\circledvee$ is well defined. Now we prove the axioms of pcm. For symmetry, assume that $x\bot y$, and hence there exist $x',y'\in X$ and $p,q\in [0,1]$ such that $x=p\cdot x'$, $y=q\cdot y'$, $p+q\leq 1$. Then $y\bot x$ via $y',x'$ and $q,p$. 
Moreover, 
\begin{align*}
    x\circledvee y & = x'+_p(y'+_{\frac{q}{1-p}}\star) \tag{\eqref{eq:defcircledvee}}\\
    & = (y'+_{\frac{q}{1-p}}\star) +_{1-p} x' \tag{symmentry in \eqref{eq:pca}}\\
    & = y'+_q(\star +_{\frac{1-p-q}{1-q}} x') \tag{associativity in \eqref{eq:pca}}\\
    & = y'+_q(x'+_{\frac{p}{1-q}}\star) \tag{symmetry in \eqref{eq:pca}}\\
& = y\circledvee x \tag{\eqref{eq:defcircledvee}}
\end{align*}
To prove that $0\bot x$ and $0\circledvee x =x$, we observe that $0=\star = 0\cdot \star$ and $x=1\cdot x$, and hence by \eqref{eq:defbot}, $0\bot x$. Obviously, $0\circledvee x = x +_1 \star = x$. Associativity is the only point in which we use the freeness of the pca. Recall that the free pca over a set $S$ is given by the set of finitely supported probability subdistributions over $S$, i.e., functions $d\colon S \to [0,1]$ such that the set $\{s\in S \mid d(s)\neq 0\}$ is finite and $\sum_{s\in S} d(s) \leq 1$. Hence, assume that $X$ is the free pca generated by a set $S$. Given $x,y,z\in X$ such that $x\bot y$ via $x',y'$ and $p,q\in [0,1]$ and $(x\circledvee y) \bot z$ via $w,z'$ and $t,u\in [0,1]$, we have to prove that $ y\bot z$ and $x \bot (y\circledvee z)$ and that $(x\circledvee y) \circledvee z = x\circledvee (y \circledvee z)$. Hence, we have that the sum of the subdistributions $p\cdot x'$, $q\cdot y'$ and $u\cdot z'$ is a subdistribution $(p\cdot x' + q\cdot y' ) + u\cdot z'= t\cdot w + u\cdot z'$. In particular, we have that $p\cdot\sum_{s\in S} x'(s) + q\cdot \sum_{s\in S} y'(s) + u\cdot \sum_{s\in S} z'(s) \leq 1$. Then, we have $y\bot z$ via $y'',z'$ and $q',u'$ where 
\[ q'\defeq q\cdot \sum_{s\in S} y'(s) \quad u' \defeq u\cdot \sum_{s\in S} z'(s)\]
and
\[y''\defeq \frac{y'}{\sum_{s\in S} y'(s)} \quad z''\defeq \frac{z'}{\sum_{s\in S} z'(s)}\text.\]
Indeed, we have that $q'+u' \leq 1$ and that $y=q'\cdot y''$ and $z=u'\cdot z''$. Similarly, we have that $x \bot (y\circledvee z)$ via $x'',w'$ and $p'',t'$, where
\[ p''\defeq p\cdot \sum_{s\in S} x'(s) \quad t' \defeq q\cdot \sum_{s\in S} y'(s) + u\cdot \sum_{s\in S} z'(s)\]
and
\[x''\defeq \frac{x'}{\sum_{s\in S} x'(s)} \quad w'\defeq \frac{q'}{t'}\cdot y'' + {\frac{u'}{t'}}\cdot z''\text.\]
Indeed, we have that $p''+t' \leq 1$ and that $x=p''\cdot x''$ and that $y\circledvee z = t' \cdot w'$. Finally, the equality $(x\circledvee y) \circledvee z = x\circledvee (y \circledvee z)$ is obvious.
\end{proof}

\begin{proposition}\label{prop:pca-pcm enrichment}
Any free pca-enriched category is a pcm-enriched category with the enrichment given by the pcm described above.
\end{proposition}
\begin{proof}
Let $\Cat{C}$ be a free pca-enriched category. By Lemma \ref{lemma:pca-pcm}, any homset $\Cat{C}[X,Y]$ is a pcm. To conclude that it is $\Cat{C}$ is pcm-enriched is enough to prove  for all properly typed arrows
\begin{itemize}
\item if $g\bot h$, then $f;(g \circledvee h) = (f;g) \circledvee (h;h) $;
\item if $f\bot g$, then  $(f \circledvee g); h = (f;h) \circledvee (g;h)$;
\item $f;0=0=0;f$.
\end{itemize}
We illustrate below the proof of $(f \circledvee g); h = (f;h) \circledvee (g;h)$.

Assume that $f\bot g$, then by \eqref{eq:defbot}, there exist arrows $f',g'$ and  $p,q\in [0,1]$ such that
\[f=p\cdot f' \quad g=q\cdot g' \qquad p+q\leq 1\text{.}\]
%Thus, by \eqref{}, \marginpar{INSERT DEFINITION ABOVE} $f\circledvee g = f' +_{p} (g'+_{\frac{q}{1-p}} \star)$. 
Thus 
\begin{align*}
(f\circledvee g); h & = (f' +_{p} (g'+_{\frac{q}{1-p}} \star));h \tag{\eqref{eq:defcircledvee}} \\
&  = (f';h +_{p} (g';h +_{\frac{q}{1-p}} \star)) \tag{$\Cat{C}$ is pca-enriched} \\
&= f;h \circledvee g;h
\end{align*}
where the last step holds by \eqref{eq:defcircledvee} and the fact that,  by Lemma \ref{lemma:initialproperties2}.\ref{lemma:p;}, $p\cdot (f';h) = (p\cdot f'); h = f;h$ and $q\cdot (g';h) = (q\cdot g'); h = g;h$.
\end{proof}

\begin{proposition}%\label{prop:fpac}
Let $\Cat{C}$ be a free pca-enriched category. If $\Cat{C}$, with the pcm-enrichment induced by Proposition~\ref{prop:pca-pcm enrichment}, is a finitely partially additive category with zero object, then it is a convex biproduct category.
\end{proposition}
\begin{proof}
Since $\Cat{C}$ is a finitely partially additive category with zero object, then by definition (\cite[Definition 3.1.2]{chophd} ) it has finite coproducts. Thus, it has zero object, i.e., initial and final object coincide, and the equality in \eqref{eq:delta} holds.

Now we show that $X_1\oplus X_2$, with the projections $\pi_1\colon X_1\oplus X_2 \to X_1$ and $\pi_2\colon X_1\oplus X_2 \to X_2$, given respectively by $(\id{X_1}\oplus\, \bangp{X_2})$ and $(\,\bangp{X_1}\oplus \id{X_2})$ is a convex product in $\Cat{C}$. Take two arrows $f\colon A\to X_1$ and $g\colon A \to X_2$ and $p_1,p_2\in (0,1)$ such that $p_1+p_2\leq 1$.
Consider now the arrows $(p_1\cdot f) ; \iota_1$ and $(p_2 \cdot g) ;\iota_2$ in the pcm  $\Cat{C}[A , X_1\oplus X_2]$. Since by Lemma \ref{lemma:initialproperties2}.\ref{lemma:p;}  $(p_1\cdot f) ; \iota_1 = p_1 \cdot(f;\iota_1)$ and 
$(p_2\cdot g) ; \iota_2 = p_2 \cdot(g;\iota_2)$ and since $p_1+p_2\leq 1$, then by definition of $\bot$ in \eqref{eq:defbot}, it holds that 
\[(p_1\cdot f) ; \iota_1 \; \bot \; (p_2 \cdot g ) ;\iota_2 \text{.}\]
We thus can take $\langle f,g\rangle_{p_1,p_2}$ as $((p_1\cdot f) ; \iota_1)  \circledvee ((p_2 \cdot g ) ;\iota_2)\colon A \to X_1\oplus X_2$ which, by  \cite[Proposition 3.1.9]{chophd}, is the unique arrow such that \[\langle f,g\rangle_{p_1,p_2} ; \pi_1 = p_1\cdot f \quad \text{ and } \quad \langle f,g\rangle_{p_1,p_2} ; \pi_2 = p_2\cdot g\text{.}\]
Condition $\ref{eq:delta}$ is trivially satisfied.
%The compatibility condition in \eqref{eq: assioma aggiuntivo} follows immediately from \cite[Proposition 3.1.8]{chophd}.
\end{proof}

\iffalse
\marginpar{Questa proposizione va rivista, smepre per via della decomposizione!}


%%%%%%%%%%%%%%%%%%%%%%%%%LA PROP NON VALE, CI SONO ANCHE QUI MAPPE COMPATIBILI NON SOMMABILI%%%%%%%%%%%%%%%%%%%%%%%%%%%%


\begin{proposition}\label{prop:T(C) is a FinPAc}
    Any free convex biproduct category $\CatT{C}$ is a finitely partially additive category with the pcm-enrichment induced by Proposition~\ref{prop:pca-pcm enrichment}.
\end{proposition}
\begin{proof}
Since $\CatT{C}$ is a convex biproduct category, then it has finite coproducts. Moreover, the pca-enrichement is freely generated, hence by Proposition~\ref{prop:pca-pcm enrichment}, it is pcm-enriched. Lemma~\ref{decomposition} and Lemma~\ref{cancellativity} imply that assumptions in Proposition~\ref{prop:B} are satisfied, hence it has jointly monic projections. Now, we prove compatible axiom of \cite[Definition 3.1.2]{chophd}: if two arrows $f,g\colon A \to X$
are compatible, which means that there exists some $h\colon A\to X \oplus X$ such that $h;\pi_i=f_i$, then they are summable. Lemma~\ref{decomposition} implies that $h$ is of the form $\diagvecp{X};(f'\oplus  g')$ for some $f',g'\colon A \to X$ and $\vec{p}=(p,q)$ with $p+q\leq 1$. Thus, $h;\pi_1 = p\cdot f'=f$ and $h;\pi_2 = q\cdot g' = g$. Hence, by \eqref{eq:defbot}, $f\bot g$ and we can define $f\circledvee g$ as in \eqref{eq:defcircledvee}. The untying axiom of \cite[Definition 3.1.2]{chophd} holds trivially.
 %if $f,g\colon A \to X$ are summable, i.e., $f\bot g$, then $(f;\iota_1) \bot (g;\iota_2)$. Summability means that there exist $f',g'\colon A \to X$ and $p,q\in [0,1]$ such that $f=p\cdot f'$, $g=q\cdot g'$ and $p+q\leq 1$. Thus,  $(f;\iota_1) \bot (g;\iota_2)$ via $f';\iota_1$, $g';\iota_2$ and $p,q$.
\end{proof}
\fi


\begin{remark}%\label{rmk:eff}
The reader may now wonder whether any convex biproduct category is also a finitely partially additive category. This is not the case in general. 

The definition of finitely partially additive category requires that \emph{compatible maps}, i.e., those $f_1\colon A \to X$, $f_2\colon A \to X$ for which there exist some $h\colon A\to X \oplus X$ such that $h;\pi_i=f_i$ (see \cite[Definition 3.1.1]{chophd} for more details), are summable. In a convex biproduct category there might exist compatible $f_1$ and $f_2$ that are not summable,  namely there are no $f_1',f_2'$ such that $f_1=p_1\cdot f_1'$ and $f_2=p_2\cdot f_2'$ for some scalar $p_1+p_2 \leq 1$. 

The essential difference can be understood as follow: summability in convex biproduct categories is fixed a-priori by the structure of pca, while in partially additive category is left to the structural properties of morphisms. For instance, while in convex biproduct categories $\oplus$ behaves as  convex product, in partially additive category $\oplus$ acts as a \emph{partial} biproduct: the product property only holds for compatible maps.   

This difference is crucial for our approach which is grounded on monoidal algebra. Indeed, in an arbitrary finitely partially additive category, one cannot define $\diagp{X}\colon X \to X\oplus X$ unless some notion of scalar is explicitly required. 
\marginpar{togliere questo ultimo paragrafo e aggiungere discorso TC non finpac}
Scalars appears in \emph{effectuses} \cite{introductioneffectus} as arrows of type $1 \to 1$ where $1$ is a distinguished object. Interestingly in $\CatT{\Cat{C}}$, for some monoidal category $(\Cat{C}, \otimes, 1)$,
every  $p\in [0,1]$ defines the arrow $\diagp{1}; (\id{1} \oplus \, \bang{1}) \colon 1 \to 1$. Nevertheless, not all the arrows in $\CatT{\Cat{C}}[1,1]$ are $p\in[0,1]$ (unless $\Cat{C}[1,1]$ is the singleton set): see for instance the syntactic categorty $\Cat{T}(\Diag{B})$ in Section \ref{sec:probboolcircuits}. To make the correspondence with effectuses more precise, one needs to consider more structure for tapes, such as their \emph{copy-discard} variants introduced in \cite{bonchi2025tapediagramsmonoidalmonads}. 
\end{remark}
\end{comment}